\section{ Gamma integrals}\label{sect:two}

\subsection{Definitions and basic properties}\label{subsec:def}

Let $u$, $\bar u$ be a pair of complex numbers of the form
\begin{gather}\label{def:z}
u=\frac n2+\nu, \qquad \bar u=-\frac n2+\nu,
\end{gather}
where $n$ is an integer and $\nu$ is complex number. We will use the notations $[u] = u-\bar u=n$ for the discrete
part and $\langle u \rangle = \nu$ for the continuous part and put $\|u\|^2=-u\bar u= -\nu^2+n^2/4$ so that for
imaginary $\nu$, $\|u\|^2\geq 0$. The $\boldsymbol \Gamma$~function of the complex field $\mathbb C$~\cite{GelfandGraevRetakh04} is defined as
\begin{gather*}
{\boldsymbol{\Gamma}}(u,\bar u)=\frac{\Gamma(u)}{\Gamma(1-\bar u)}
=\frac{\Gamma(n/2+\nu)}{\Gamma(1+n/2-\nu)}=(-1)^n\frac{\Gamma(-n/2+\nu)}{\Gamma(1-n/2-\nu)} =(-1)^{[u]}\Gamma(\bar u, u).
\end{gather*}
In what follows we will, for brevity, display only the first argument of the $\boldsymbol \Gamma$ function, i.e.,
$\boldsymbol \Gamma(u)\equiv \boldsymbol \Gamma(u,\bar u)$. Hereafter the following functional relations will be useful
\begin{gather}\label{func:gamma}
\boldsymbol \Gamma(u)\boldsymbol \Gamma(1-u)=(-1)^{[u]}, \qquad
\boldsymbol \Gamma(u+1)=-u\bar u \boldsymbol \Gamma(u).
\end{gather}
The $\boldsymbol \Gamma$ function appears in the generalization of Gustafson's integrals to the complex case.

The corresponding integrals take the following form~\cite{Derkachov:2016ucn,Derkachov:2017pvx}
\begin{subequations}\label{Gustafson12}
\begin{gather}
\frac1{N!} \sum_{n_1,\dots,n_N\in\mathbb{Z}+\frac{\sigma}{2}}
\int_{-{\rm i}\infty}^{{\rm i}\infty}
\frac{\prod\limits_{m=1}^{N+1}\prod\limits_{k=1}^{N}
\boldsymbol{\Gamma}(z_m-u_k)\boldsymbol{\Gamma}(u_k+w_m)}{ \prod\limits_{m<j}
\boldsymbol{\Gamma}(u_m- u_j)\boldsymbol{\Gamma}(u_j- u_m)}
\frac{{\rm d}\nu_1}{2\pi {\rm i}}\cdots\frac{{\rm d}\nu_N}{2\pi {\rm i}}\nonumber\\
\qquad{} =\frac{\prod\limits_{k,j=1}^{N+1}\boldsymbol{\Gamma}(z_k+ w_j)}{
\boldsymbol{\Gamma}\left(\sum\limits_{k=1}^{N+1} (z_k+w_k) \right)},\label{gustafson1}\\
\frac1{2^{N}N!}\sum_{n_1,\dots,n_N\in\mathbb{Z}+\frac{\sigma}{2}}
\int_{-{\rm i}\infty}^{{\rm i}\infty}
\frac{ \prod\limits_{k=1}^{N} \prod\limits_{m=1}^{2N+2}
 \boldsymbol{\Gamma}(z_m\pm u_k)}
{\prod\limits_{k=1}^{N}{\boldsymbol{\Gamma}}(\pm 2u_k) \prod\limits_{ k<j} {\boldsymbol{\Gamma}} (\pm u_k\pm u_j)
}\frac{{\rm d}\nu_1}{2\pi {\rm i}}\cdots\frac{{\rm d}\nu_N}{2\pi {\rm i}}\nonumber\\
\qquad{} = \frac{\prod\limits_{j<k}{\boldsymbol{\Gamma}}(z_j + z_k)}{{\boldsymbol{\Gamma}}\left(\sum\limits_{k=1}^{2N+2} z_k\right)}, \label{gustafson2}
\end{gather}
\end{subequations}
where we put $\boldsymbol \Gamma(a\pm b)\equiv \boldsymbol \Gamma(a+b) \boldsymbol \Gamma(a-b)$, $\boldsymbol
\Gamma(\pm a\pm b)\equiv \boldsymbol \Gamma(a+b) \boldsymbol \Gamma(a-b)\Gamma(-a+b) \boldsymbol \Gamma(-a-b)$. We will
refer to the integrals in the first and second lines as $I_N^{(1)}$ and $I_N^{(2)}$, respectively.

The variables $u_k$, $w_m$, $z_m$ have the form~\eqref{def:z}
\begin{gather*}
u_r=\frac{n_r}2+ \nu_r, \qquad z_j=\frac{m_j}2+ x_j, \qquad w_m=\frac{\ell_m}2+ y_m
\end{gather*}
and similarly for the barred variables. However, $n_r$, $m_i$, $\ell_m$ are allowed take integer or half-integer values,
simultaneously. Accordingly, the sums in~\eqref{Gustafson12} go over integers ($\sigma=0$) or half-integers
($\sigma=1/2)$. For the first integral~\eqref{gustafson1} there is no difference between the integer/half-integer cases
since they are related by the change of variables: $[u_r], [w_m], [z_j] \mapsto [u_r]+1/2, [w_m]-1/2,[z_j]+1/2$ so that we will assume that the variables $[u_k]$, $[w_m]$, $[z_m]$ in the integral~\eqref{gustafson1} are integers.

The integration contours in \eqref{Gustafson12} separate the series of $``\pm"$ poles due to the $\boldsymbol \Gamma$ functions in the numerators. The poles are located at
\begin{gather}\label{poles-I}
\nu^{\text{I},+}_{rj}(p)=\frac12{|n_r-m_j|} +x_j +p, \qquad \nu^{\text{I},-}_{rj}(p)=-\frac12{|n_r+\ell_j|} -y_j -p,
\qquad p\geq 0,
\end{gather}
where $r\in\{1, N\}$, $j\in\{1,N+1\}$ for the first integral and
\begin{gather*}
\nu^{\text{II},+}_{rj}(p)=\frac12{|n_r-m_j|} +x_j +p, \qquad \nu^{\text{II},-}_{rj}(p)=-\frac12{|n_r+m_j|} -x_j -p,\qquad p\geq 0,
\end{gather*}
where $r\in\{1, N\}$, $j\in\{1,2N+2\}$, for the second one.

Let us discuss now the convergence properties of the integrals~\eqref{Gustafson12}. Since the integrands are meromorphic functions and contours of integration avoid poles
it is sufficient to analyse the region of large $u_r=n_r/2+ {\rm i}\nu_r$, $ \|u_r\|^2=\nu_r^2+n_r^2/4 \to\infty$, only.
With the help of \eqref{func:gamma} we simplify the denominators in the integrals~\eqref{Gustafson12} as follows
\begin{gather}\label{den1}
\prod_{1\leq i<k\leq N} \frac1{\boldsymbol{\Gamma}(u_i- u_k)\boldsymbol{\Gamma}(u_k- u_i)} = (-1)^{(N+1)\sum_k [u_k] } \prod_{1\leq i<k\leq N}\|u_i-u_k\|^2
\end{gather}
and
\begin{gather}
\prod_{k=1}^{N}\frac1{\boldsymbol{\Gamma}(\pm 2u_k)} \prod_{1\leq i<m\leq N}\frac1{\boldsymbol \Gamma (\pm u_i\pm u_m)}\nonumber\\
\qquad{} =\varkappa_N 4^{N} \prod_{k=1}^N \|u_k\|^2
\prod_{1\leq i<m\leq N}\|u_i-u_m\|^2\|u_i+u_m\|^2,\label{den2}
\end{gather}
where $\varkappa_N=1$ for the integer case and $\varkappa_N=(-1)^{N(N+1)/2}$ for the half-integer case. Finally, taking
into account that for large $u$
\begin{gather}
\boldsymbol \Gamma(z -u) \boldsymbol \Gamma(u+w) =
(-1)^{[z-u]}\frac{\boldsymbol{\Gamma}(u+w)}{\boldsymbol{\Gamma}(u-z)} \nonumber\\
\hphantom{\boldsymbol \Gamma(z -u) \boldsymbol \Gamma(u+w)}{} = (-1)^{[z-u]}
 u^{z+w-1} (-\bar u)^{\bar z+\bar w -1}\bigl(1+O\left(1/\|u\|\right)\bigl) \label{large-u-as}
\end{gather}
we conclude that the integrals~\eqref{Gustafson12} converge absolutely provided
\begin{gather}\label{convergence}
 \sum_{j=1}^{N+1} \operatorname{Re}(x_j+y_j) <1 \qquad \text{and} \qquad \sum_{j=1}^{2N+2} \operatorname{Re} (x_j) <1,
\end{gather}
respectively. From now on we assume that these conditions are satisfied.

\subsection{Determinant representation}\label{subsec:det}

In this subsection we present the integrals~\eqref{Gustafson12} as determinants of $1$-dimensional integrals. Such a representation will be useful in what follows.\footnote{The determinant representation for elliptic hypergeometric integrals was constructed in~\cite{RainSpiridonov}.} For its derivation let us denote by $\mathcal Q (u|z,w)$ the function
\begin{gather*}%\label{QGamma}
\mathcal Q (u|z,w)={\prod_{k=1}^{N+1}(-1)^{[u]}\boldsymbol\Gamma(z_k-u)\boldsymbol\Gamma(u+w_k)}
\end{gather*}
and by ${\mathcal{Q}}_{ik}(z,w)$ its Mellin moments
\begin{gather}\label{MellinQ}
{\mathcal{Q}}_{ik}(z,w)=\int \mathcal Du u^{i-1} (-\bar u)^{k-1} \mathcal Q (u|z,w), \qquad i,k=1,\ldots,N.
\end{gather}
Here we introduced a short-hand notation for the integration measure
\begin{gather*}
\int\mathcal Du\equiv
\sum_{n=-\infty}^{\infty}\int_{-{\rm i}\infty}^{{\rm i}\infty}\frac{{\rm d}\nu}{2\pi {\rm i}} .
\end{gather*}
Let $\boldsymbol{\mathcal Q}_N(z,w)$ be the following $N\times N$ matrix constructed from the Mellin moments
 \begin{gather}\label{Q-matrix}
\boldsymbol{\mathcal Q}_N(z,w)=\begin{pmatrix}
{\mathcal{Q}}_{11}(z,w)&\cdots& {\mathcal{Q}}_{1N}(z,w)
\\
\vdots & \ddots &\vdots
\\
{\mathcal{Q}}_{N1}(z,w)&\cdots& {\mathcal{Q}}_{NN}(z,w)
\end{pmatrix}.
\end{gather}
Rewriting the product on the r.h.s.\ of equation~\eqref{den1} as the product of two Vandermonde determinants
\begin{gather*}
\prod_{1\leq i<k\leq N}\|u_i-u_k\|^2=\Delta(u) \Delta (-\bar u), \qquad \Delta(u)=\det_{1\leq i,j\leq N} u_j^{i-1}
\end{gather*}
and taking into account the symmetry of the integrand in~\eqref{gustafson1} with respect to the permutations $u_i\leftrightarrow u_j$ we bring the first integral into the determinant form
\begin{gather}\label{det-rep-1}
I_N^{(1)} =\int\mathcal Du_1\cdots \mathcal Du_N \Delta(-\bar u) \prod_{k=1}^N \mathcal Q (u_k|z,w) u_k^{k-1}=\det \boldsymbol{\mathcal Q}_N(z,w).
\end{gather}
Proceeding in a similar way one gets the determinant representation for the second integral as follows
\begin{gather*}%\label{detQ}
I_N^{(2)} =\varkappa_N \det \boldsymbol{\widetilde {\mathcal Q}}_N(z)=\varkappa_N \det \begin{vmatrix}
\widetilde{\mathcal{Q}}_{11}(z)&\cdots& \widetilde{\mathcal{Q}}_{1N}(z)\\
\vdots & \ddots &\vdots\\
\widetilde{\mathcal{Q}}_{N1}(z)&\cdots& \widetilde{\mathcal{Q}}_{NN}(z)
\end{vmatrix},
\end{gather*}
where $\varkappa_N$ is a phase factor, see equation~\eqref{den2},
\begin{gather*}
\widetilde {\mathcal Q}_{ik}(z)=2\int \mathcal{D}u u^{2i-1}(-\bar u)^{2k-1} \widetilde {\mathcal Q}(u,z) \qquad\text{and}
\qquad
\widetilde{ \mathcal Q}(u,z)=\prod_{j=1}^{2N+2} \boldsymbol\Gamma(z_j\pm u).
\end{gather*}
Let us note here that the conditions~\eqref{convergence} are equivalent to the requirement of absolute convergence of the Mellin moments $\mathcal Q_{NN}(z,w)$ and $\widetilde{\mathcal Q}_{NN}(z)$, respectively.

\subsection{Proof of identities~(\ref{Gustafson12})}\label{subsec:proof}

Calculating the integral~\eqref{gustafson1} we will assume that the parameters $z_k$, $w_k$ satisfy the conditions
\begin{gather}\label{zw-cond}
\sum_{k=1}^{N+1} \operatorname{Re}(z_k+w_k)<1 \qquad \text{and} \qquad \sum_{k=1}^{N+1} \operatorname{Re} (\bar z_k+\bar w_k)<1.
\end{gather}
These conditions imply the condition~\eqref{convergence}, but do not follow from it and will be removed at the end of the calculation.

By virtue of~\eqref{convergence} the integrals~\eqref{MellinQ} are all absolutely convergent. We evaluate the integral
over $\nu$ by closing the contour in the left half-plane and then computing the sum over the residues. Recalling that
$u=n/2+\nu$ and $\bar u=-n/2+\nu$ one gets
\begin{gather*}
M_{ik}(n) =\int^{{\rm i}\infty}_{-{\rm i}\infty} u^{i-1}(-\bar u)^{k-1} \mathcal Q(u|z,w) \frac{{\rm d}\nu}{2\pi {\rm i}}\notag\\
\hphantom{M_{ik}(n)}{} =
(-1)^{(N+1)n}\sum_{j=1}^{N+1} \sum_{p=0}^{\infty}\frac{(-1)^p}{p! \bar p_j!} (-w_{j}-p)^{i-1} (\bar w_{j} + \bar p_j)^{k-1}\notag\\
\hphantom{M_{ik}(n)=}{} \times
\prod_{m=1}^{N+1} \frac{\Gamma(z_m+w_{j}+p)}{\Gamma(1-\bar z_m-\bar w_{j}-\bar p_j)}
\prod_{m \neq j}^{N+1} \frac{\Gamma(w_m-w_{j}-p)}{\Gamma(1-\bar w_m+ \bar w_{j}+ \bar p_j)},%\label{M-nu-n}
\end{gather*}
where $\bar p_j = p+n+\ell_j$.

The first observation is that since the summand vanishes for $\bar p_j < 0$, only the poles at
$\nu^{\mathrm{I},-}_{rj}(p)$ given in~\eqref{poles-I}, contribute to the integral and, second, that under the
assumptions~\eqref{zw-cond} the sum over $p$ converges uniformly on $n$. Therefore, to evaluate $\sum_n M_{ik}(n)$ we
swap the summation over $n$ and $p$, change the summation variable from $n$ to $\bar p_j$ and finally suppress the
index $j$ in~$\bar p_j$: $\bar p_j\mapsto \bar p$
\begin{gather*}
\sum_n M_{ik}(n)=\sum_{j=1}^{N+1}\sum_{p= 0}^{\infty} \sum_{n=-\infty}^{\infty}( \ldots )
\mapsto \sum_{j=1}^{N+1}\sum_{p= 0}^{\infty}
\sum_{\bar p_j=-\infty}^{\infty}( \ldots ) =\sum_{j=1}^{N+1}\sum_{p= 0}^{\infty}
\sum_{\bar p =0}^{\infty}( \ldots ).
\end{gather*}
The final answer can be written in the form
\begin{gather}
\mathcal Q_{ik}(z,w) =\sum_{j=1}^{N+1} (-1)^{(N+1)[w_j]}
\left(\prod_{m=1}^{N+1}\frac{\Gamma(z_m+w_{j})}{\Gamma(1-\bar z_m-\bar w_{j})}\right)\left(
\prod_{\substack{m=1\\ m\neq j}}^{ N+1}\frac{ \Gamma(w_m-w_{j})}{\Gamma(1-\bar w_m+\bar w_{j})}\right)\notag\\
\hphantom{\mathcal Q_{ik}(z,w)=}{} \times
\left(\sum_{p=0}^{\infty} (-w_{j}-p)^{i-1}\prod_{m=1}^{N+1}
\frac{(z_m+w_{j})_{p}}{(1-w_m+w_{j})_{p}}\right)\nonumber\\
\hphantom{\mathcal Q_{ik}(z,w)=}{} \times
\left(\sum_{\bar p=0}^{\infty}(\bar w_{j} + \bar p)^{k-1}\prod_{m=1}^{N+1}
\frac{(\bar z_m+\bar w_{j})_{\bar p}}{(1-\bar w_m+\bar w_{j})_{\bar p}}\right) , \label{Q-factorized}
\end{gather}
where $(a)_p$ is the Pochhammer symbol. Thus $\mathcal Q_{ik}=\sum\limits_{j=1}^{N+1} \mathcal Q_{ik}(j)= \sum\limits_{j=1}^{N+1} \mathcal Q_{i}(j)\bar{\mathcal Q}_k(j)$. Hence the determinant can be represented as follows
\begin{gather*}
\det\mathcal Q =\sum_{j_1=1}^{N+1}\cdots \sum_{j_N=1}^{N+1}
\det\|\mathcal Q_{ik}(j_k)\| =\sum_{\sigma\in S_{N+1}}
\det \| \mathcal Q_{ik}(j_{\sigma(k)}\|
\notag\\
\hphantom{\det\mathcal Q}{} =\sum_{\sigma\in S_{N+1}} \left(\prod_{k=1}^N\bar{\mathcal Q}_k(j_{\sigma(k)})\right)
\det\| \mathcal Q_{i}(j_{\sigma(k)})\|,
\end{gather*}
where we take into account that $\det\|Q_{ik}(j_k)\|=0$ whenever $j_{m}=j_n$ for $n\neq m$. Then making use of~\eqref{Q-factorized} one can bring~\eqref{det-rep-1} into the following form
\begin{gather*}
I_N^{(1)}
=\frac1{N!}\!\sum_{\sigma\in S_{N+1}}\! (-1)^{(N+1)\sum\limits_{s=1}^N
[w_{\sigma(s)}]}\prod_{k=1}^{N}\left(
\prod_{j=1}^{N+1}\frac{\Gamma(z_j+w_{\sigma(k)})}{\Gamma(1-\bar z_j-\bar w_{\sigma(k)})}
\!\prod_{\substack{m=1\\ m\neq \sigma(k)}}^{N+1}\!
 \frac{ \Gamma(w_m-w_{\sigma(k)})}{\Gamma(1-\bar w_m+\bar w_{\sigma(k)})}\!\right)
\\
\hphantom{I_N^{(1)} =}{}
\times
\sum_{p_1,\ldots,p_N= 0}^\infty
 \left(\prod_{k<m} (w_{\sigma(k)}+p_k-w_{\sigma(m)}-p_m)\right)\prod_{m=1}^{N+1} \prod_{k=1}^{N}
\frac{(z_m+w_{\sigma(k)})_{p_k}}{(1-w_m+w_{\sigma(k)})_{p_k}}
\notag\\
\hphantom{I_N^{(1)} =}{}
\times
\sum_{\bar p_1,\ldots,\bar p_N=0}^\infty
 \left(\prod_{k<m} (\bar w_{\sigma(m)}+\bar p_m-\bar w_{\sigma(k)}-\bar p_k)\right)
\prod_{m=1}^{N+1} \prod_{k=1}^{N}
\frac{(\bar z_m+\bar w_{\sigma(k)})_{\bar p_k}}{(1-\bar w_m+\bar w_{\sigma(k)})_{\bar p_k}} .
\end{gather*}
The infinite sums over $\{ p\}$, $\{\bar p\}$ can be evaluated by Milne's ${\rm U}(n)$ Gauss
summation~\cite{MILNE198859}, see also~\cite[equation~(5.8)]{Gustafson94} and~\cite{Schlosser},
\begin{gather}
\sum_{p_1,\ldots,p_N = 0}^\infty \left(\prod_{1\leq k<m\leq N} (\alpha_{\sigma(k)}+p_k-\alpha_{\sigma(m)}-p_m)\right)
\prod_{i=1}^{N+1} \prod_{k=1}^{N}\frac{(\beta_i+\alpha_{\sigma(k)})_{p_k}}{(1-\alpha_i+\alpha_{\sigma(k)})_{p_k}}\nonumber\\
 \qquad{}=\frac{\Gamma\left(1-\sum\limits_{k=1}^{N+1}(\alpha_k+\beta_k)\right)
\prod\limits_{i=1}^{N}\Gamma(1+\alpha_{\sigma(i)}-\alpha_{\sigma(N+1)})}
{\prod\limits_{i=1}^{N+1}\Gamma(1-\beta_i-\alpha_{\sigma(N+1)})}.\label{Milne}
\end{gather}
Using~\eqref{Milne} we obtain the following representation for the integral $I_N^{(1)}$
\begin{gather*}
I_N^{(1)} =\frac1{N!}\!\sum_{\sigma\in S_{N+1}} \!(-1)^{(N+1)\sum\limits_{s=1}^N
[w_{\sigma(s)}] }\prod_{k=1}^{N}\left(
\prod_{i=1}^{N+1}\frac{\Gamma(z_i+w_{\sigma(k)})}{\Gamma(1-\bar z_i-\bar w_{\sigma(k)})}
\!\prod_{\substack{m=1\\ m\neq \sigma(k)}}^{N+1}\!
 \frac{ \Gamma(w_m-w_{\sigma(k)})}{\Gamma(1-\bar w_m+\bar w_{\sigma(k)})}\!\right)\\
\hphantom{I_N^{(1)} =}{} \times
\prod_{1\leq m<j\leq N} (w_{\sigma(m)}-w_{\sigma(j)})
\Gamma\left(1-\sum_{k=1}^{N+1}(z_k+w_k)\right)
\frac{\prod\limits_{k=1}^{N}\Gamma(1+w_{\sigma(k)}-w_{\sigma(N+1)})}
{\prod\limits_{k=1}^{N+1} \Gamma(1-z_k-w_{\sigma(N+1)}))}\\
\hphantom{I_N^{(1)} =}{} \times \prod_{1\leq m<j\leq N}
(\bar w_{\sigma(j)}-\bar w_{\sigma(m)})
\Gamma\left(1-\sum_{k=1}^{N+1}(\bar z_k+\bar w_k)\right)
\frac{\prod\limits_{k=1}^{N}
\Gamma(1+\bar w_{\sigma(k)}-\bar w_{\sigma(N+1)})}{
\prod\limits_{k=1}^{N+1}
\Gamma(1-\bar z_k-\bar w_{\sigma(N+1)})}.
\end{gather*}
After some simplifications this can be written as
\begin{gather*}
I_N^{(1)}=\frac{\Gamma\left(1-\sum\limits_{k=1}^{N+1}(\bar z_k+\bar w_k)\right)}{\Gamma\left(\sum\limits_{k=1}^{N+1}( z_k+ w_k)\right)}
\prod_{i,k=1}^{N+1}\frac{\Gamma(z_i+w_k)}{\Gamma(1-\bar z_i-\bar w_k)}
\frac{ R_N}{N!\sin\pi\left(\sum\limits_{k=1}^{N+1}( z_k+ w_k)\right)},
\end{gather*}
where
\begin{gather}
R_N = \sum_{\sigma\in S_{N+1}}
\frac{\prod\limits_{k=1}^{N+1}
\sin\pi(z_{\sigma(k)}+w_{\sigma(N+1)})}
{\prod\limits_{k=1}^N\sin\pi(w_{\sigma(N+1)}-w_{\sigma(k)})}
\nonumber\\
\hphantom{R_N =}{} \times(-1)^{(N+1)\sum\limits_{s=1}^N[w_{\sigma(s)}]}\prod_{1\leq k<j\leq N}
\frac{\sin\pi(\bar w_{\sigma(j)}-\bar w_{\sigma(k)})}{\sin\pi( w_{\sigma(j)}- w_{\sigma(k)})} .\label{R_N}
\end{gather}
Taking into account that $w_k-\bar w_k$ is an integer one finds that the last product in~\eqref{R_N} yields $(-1)^{(N-1)\sum\limits_{s=1}^N [w_{\sigma(s)}]}$ which cancels the second factor in~\eqref{R_N}. In the last step we use the \cite[Lemma~5.10]{Gustafson94} which states that
\begin{gather*}
\sum_{\sigma\in S_{N+1}}
\frac{\prod\limits_{k=1}^{N+1}
\sin\pi(\beta_{k}+\alpha_{\sigma(N+1)})}
{\prod\limits_{k=1}^N\sin\pi(\alpha_{\sigma(N+1)}-\alpha_{\sigma(k)})} =
N!\sin\pi \sum_{k=1}^{N+1}(\alpha_k+\beta_k).
\end{gather*}
It results in
\begin{gather*}
R_N=N!\sin\pi \sum_{k=1}^{N+1}( z_k+ w_k),
\end{gather*}
so that we get the required result for $I_N^{(1)}$
 \begin{gather*}
I_N^{(1)}=\frac{\prod\limits_{k,j=1}^{N+1}\boldsymbol{\Gamma}(z_k+ w_j)}{
\boldsymbol{\Gamma}\left(\sum\limits_{k=1}^{N+1} (z_k+w_k) \right)} .
\end{gather*}
